\begin{titlepage}
\begin{tikzpicture}[remember picture,overlay]
    \fill[Ink] (current page.north west) rectangle ([xshift=12mm]current page.south west);
    \fill[Accent] ([xshift=12mm]current page.north west)
        rectangle ([xshift=15mm]current page.south west);
    \fill[Paper] ([xshift=-85mm]current page.north east)
        rectangle ([yshift=-122mm]current page.north east);
    \foreach \i in {0,...,6}{
        \draw[Rule,line width=0.4pt]
            ([xshift=-78mm,yshift=-28mm-\i*12mm]current page.north east)
            -- ([xshift=-12mm,yshift=-28mm-\i*12mm]current page.north east);
    }
    \draw[Copper,line width=2pt]
        ([xshift=-78mm,yshift=-43mm]current page.north east)
        -- ([xshift=-59mm,yshift=-64mm]current page.north east)
        -- ([xshift=-40mm,yshift=-52mm]current page.north east)
        -- ([xshift=-21mm,yshift=-88mm]current page.north east);
\end{tikzpicture}
\vspace*{17mm}
{\sffamily\bfseries\color{Accent}\large MTY\enspace /\enspace NOTEBOOK}\par
\vspace{42mm}
{\fontsize{34}{44}\selectfont\sffamily\bfseries\BookTitle}\par
\vspace{7mm}
{\large\sffamily\color{Muted}\BookSubtitle}\par
\vspace{16mm}
{\color{Accent}\rule{35mm}{1.4pt}}\par
\vspace{6mm}
{\large 保留熟悉的码风，积累可复用的算法。}\par
\vspace{4mm}
{\color{Muted}说明、源码与排版分离\quad /\quad 持续维护\quad /\quad 按需打印}\par
\vfill
\begin{tabularx}{\linewidth}{@{}Xr@{}}
    {\sffamily\bfseries\BookAuthor} & {\sffamily\BookVersion} \\
    {\small\sffamily\color{Muted}\BookRepository} & {\small\sffamily\color{Muted}\BookDate}
\end{tabularx}
\vspace{4mm}
\end{titlepage}
